\documentclass[10pt,a4paper]{amsart}
\usepackage[margin=19mm]{geometry}
\usepackage{amsmath,amssymb,mathtools}
\usepackage[scr=rsfs,cal=euler]{mathalfa}
\usepackage{xfrac}
\usepackage[unicode,hidelinks,bookmarks=false]{hyperref}
\newcommand{\ie}{i.e.\ }
\newcommand{\cf}{cf.\ }
\newcommand{\resp}{respectively\ }
\newcommand{\Subsection}{\subsection}
\providecommand{\nobreakdash}{\nobreak\hbox{-}}
\setlength{\emergencystretch}{2em}
\multlinegap0pt
\usepackage{amssymb}
\usepackage{bm}
\usepackage[shortlabels]{enumitem}
\usepackage{xcolor}
\newtheorem{assumption}{Assumption}
\newtheorem{proposition}{Proposition}
\newtheorem{remark}{Remark}
\newtheorem{lemma}{Lemma}
\newtheorem{theorem}{Theorem}
\newcommand{\II}{{\mathbb I}}
\newcommand{\RR}{{\mathbb R}} % the set of real numbers
\newcommand{\abs}[1]{\left\lvert#1\right\rvert}
\newcommand{\cF}{{\mathcal F}}
\newcommand{\cP}{{\mathcal P}}
\newcommand{\cX}{{\mathcal X}}
\newcommand{\cZ}{{\mathcal Z}}
\newcommand{\ccr}{{\mathfrak r}}
\newcommand{\norm}[1]{\left\lVert#1\right\rVert}
\newcommand{\tpose}[1]{#1^{\top}}
\title{Sparsity Regularized and Robust Mean Variance Portfolio Selection Under Ellipsoidal Uncertainty}
\author{Deniz Akkaya, Emre Can Yayla, Buse \c{S}en, Mustafa \c{C}. P{\i}nar}
\date{Source: arXiv:2609.11749v1 (CC BY 4.0)}
\begin{document}
\maketitle

\noindent\emph{Source and licence.} Sparsity Regularized and Robust Mean Variance Portfolio Selection Under Ellipsoidal Uncertainty. Version arXiv:2609.11749v1 (CC BY 4.0), distributed by arXiv under the Creative Commons Attribution 4.0 International licence, http://creativecommons.org/licenses/by/4.0/, as recorded in the arXiv licence metadata for this exact version and shown on the abstract page https://arxiv.org/abs/2609.11749v1. Full licence notice: \texttt{licenses/arxiv-2609.11749-CC-BY-4.0.txt}.\par\smallskip

\noindent\textit{Editorial scope.} Counted unit: the article's theorem thm:lower-bound (source0.tex lines 407--416). The article's complete original proof of it is delivered with it (source0.tex lines 417--482, one \texttt{proof} environment of 66 lines). No mathematics is written, repaired or re-derived: the statement and the proof are the article's own lines, in the article's own notation, and no retained line is substituted or re-flowed.  Retained uncounted as the source-local context the proof needs: lines 190--195; lines 238--244; lines 261--263; lines 271--273; lines 321--324; lines 350--353; lines 379--381; lines 385--389.  Nothing was omitted from any retained span, so the package carries no omission record.  No retained line contains a citation, so the wrapper carries no bibliography and no bibliography entry.  No retained line contains a figure environment, an image-inclusion command or drawing code, so no image is delivered and the package declares no asset dependency.  Editorial formatting only, in addition: an AMS \texttt{amsart} preamble that restates the article's own declarations and macros used by the retained lines, namely the environments \texttt{assumption}, \texttt{proposition}, \texttt{remark}, \texttt{lemma}, \texttt{theorem} on separate counters, together with the standard packages those lines need.  The retained environments print in this document as Assumption 1, Proposition 1, Remark 1, Lemma 1, Theorem 1 in order of appearance, whereas the article prints the same items with the numbering of its own full document.  The complete native source of the article as distributed in the arXiv e-print for this exact version is bundled unchanged as \texttt{sources/210047/source0.tex}.
\begin{align}
\tag{$\cP^1$}\label{prob:cprmvp}
\min_{x\in\RR^N} \; \cF_\beta(x):= \tpose{x} D x + \beta \norm{x}_0
\quad \text{s.t.} \quad
\tpose{\ccr} x - \gamma \norm{x}_D \ge \bar{\ccr}.
\end{align}

\begin{assumption}\label{assump:feas}
The uncertainty radius $\gamma$ satisfies
\begin{align*}
\min_{i\in\II_N}\frac{\abs{\ccr[i]}}{\sqrt{d_i[i]}}>\gamma,
\end{align*}
where $\sqrt{d_i[i]}$ denotes the square root of $i^{th}$ diagonal entry of $D$.
\end{assumption}

\begin{align}\tag{$\cP^1_\omega$}\label{prob:mvp-omega}
	\min_{x\in \cX\cap K_\omega}\tpose{x}Dx.
\end{align}

\begin{align}\tag{$\cZ\cP^1_\omega$}\label{prob:zpmvp}
	\min_{u\in \cX_\omega}\tpose{u}D_\omega u, \quad \abs{\omega}\geq 1.
\end{align}

\begin{proposition}\label{prop:dual_solution}
For nonempty $\omega\subseteq \II_N$, the optimal Lagrange multiplier associated with the constraint of~\eqref{prob:zpmvp} is $\mu^* = \frac{2\bar{\ccr}}{(H_\omega - \gamma)^2}$.
Moreover, the constraint is active at the optimal solution $\hat u$, \textit{i.e.}, $\tpose{\ccr_\omega}\hat{u}-\gamma \norm{\hat{u}}_{D_\omega}= \bar{\ccr}$.
\end{proposition}

\begin{remark}
\label{rem:unique}
For $\omega\subseteq\II_N$ with $\abs{\omega}\geq 1$, the point $\Xi(\omega)\in\RR^N$ is the unique solution of \eqref{prob:mvp-omega}.
\end{remark}

\begin{lemma}\label{lemma8}
Let $\beta>0$ and let $\hat{x}$ be a (local) minimizer of \eqref{prob:cprmvp}. Then $\hat{x}=\Xi(\sigma(\hat{x}))$.
\end{lemma}

\begin{proposition}\label{prop:globup}
   Let $\beta>0$ and suppose $\hat{x}$ is a global minimizer of~\eqref{prob:cprmvp}. Then we have
    \begin{align*}
        \norm{\hat{x}}_\infty\leq \sqrt{\frac{\eta}{s_N(D)}}<\infty,  \;  \text{ where } \eta=\min_{i\in\II_N}\left\{\frac{\bar{\ccr}^2d_i[i]}{\left(|\ccr[i]|-\gamma\sqrt{d_i[i]}\right)^2}\right\}.\end{align*}
\end{proposition}

\begin{theorem}
\label{thm:lower-bound}
    Let $\beta>0$ and suppose $\hat{x}$ is a global minimizer of~\eqref{prob:cprmvp}. Let also $\hat{\sigma}=\sigma(\hat{x})$ denote the support of $\hat{x}$, and we define for all $i\in\II_N$
    \begin{align*}
        \rho_i=\max_{\substack{s\in\{-1,1\}\\ j\in \II_N\setminus\{i\}}}\left\{\norm{e_i+s\frac{\abs{\ccr[i]}+\gamma\sqrt{d_i[i]}}{\abs{\ccr[j]}-\gamma\sqrt{d_j[j]}}e_j}_D\right\} \quad\text{ and } \quad\eta=\min_{i\in\II_N}\left\{\frac{\bar{\ccr}^2d_i[i]}{(\abs{\ccr[i]}-\gamma\sqrt{d_i[i]})^2}\right\} .
    \end{align*} Then for every $i\in\hat{\sigma}$, we have
    \begin{align*}
    \abs{\hat{x}[i]}\geq \min\left\{\frac{\sqrt{\eta+\beta}-\sqrt{\eta}}{\rho_i},\frac{\bar{\ccr}}{\abs{\ccr[i]}-\gamma\sqrt{d_i[i]}}\right\}.
\end{align*}
\end{theorem}

\begin{proof}
    Let $i\in\hat{\sigma}$. We consider two cases based on the cardinality of $\hat{\sigma}$.

    \textbf{Case 1:} Assume $\abs{\hat{\sigma}}\geq 2$, then there exists $j\in\hat{\sigma}$ such that $i\neq j$. We define $g_{ij}:\RR^N\rightarrow \RR^N$ as $g_{ij}(x)=x-x[i]e_i-x[j]e_j$ where $e_i$ and $e_j$ are canonical basis vectors of the specified index. We introduce the function $f(t_i,t_j)=\cF_\beta(g_{ij}(\hat{x})+t_ie_i+t_je_j)$. Finally, we define a feasibility function
    \begin{align*}
        h(t_i,t_j)=\tpose{\ccr}g_{ij}(\hat{x})+t_i\ccr[i]+t_j\ccr[j]-\gamma\norm{g_{ij}(\hat{x})+t_ie_i+t_je_j}_D-\bar{\ccr}.
    \end{align*}
   Since $\hat{x}$ is a global minimizer, Lemma~\ref{lemma8} gives $\hat{x}=\Xi(\hat{\sigma})$.  By Remark~\ref{rem:unique}, $\hat{x}$ is therefore the unique solution of \eqref{prob:mvp-omega} with $\omega=\hat{\sigma}$. Proposition~\ref{prop:dual_solution} then implies that the restricted constraint is active at $\hat{x}$. This, in turn,  implies that the main constraint is also active, yielding $h(\hat{x}[i],\hat{x}[j])= 0$. Precisely, we rewrite this equality as
    \begin{align*}
        h(\hat{x}[i],\hat{x}[j])=\tpose{\ccr}\hat{x}-\gamma\norm{\hat{x}}_D-\bar{\ccr}=0\Rightarrow \bar{\ccr}=\tpose{\ccr}\hat{x}-\gamma\norm{\hat{x}}_D.
    \end{align*}
    Now, we define a map by fixing a choice for $t_j$
    \begin{align*}
        h(t_i,t_j)&=\tpose{\ccr}g_{ij}(\hat{x})+t_i\ccr[i]+t_j\ccr[j]-\gamma\norm{g_{ij}(\hat{x})+t_ie_i+t_je_j}_D-\bar{\ccr}\\
        &=\gamma\left(\norm{\hat{x}}_D-\norm{g_{ij}(\hat{x})+t_ie_i+t_je_j}_D\right)+(t_i-\hat{x}[i])\ccr[i]+(t_j-\hat{x}[j])\ccr[j]\\&
        \geq-\gamma\left(\norm{(t_i-\hat{x}[i])e_i+(t_j-\hat{x}[j])e_j}_D\right)+(t_i-\hat{x}[i])\ccr[i]+(t_j-\hat{x}[j])\ccr[j]\\&
        \geq -\gamma\abs{t_i-\hat{x}[i]}\sqrt{d_i[i]}-\gamma\abs{t_j-\hat{x}[j]}\sqrt{d_j[j]}+(t_i-\hat{x}[i])\ccr[i]+(t_j-\hat{x}[j])\ccr[j] ,
    \end{align*}
    where the first inequality exploits  $\|u\|-\|v\|\geq -\|u-v\|$ with $u=\hat x$ and $v=g_{ij}(\hat x)+ t_i e_i+t_je_j$, and the second follows from the triangle inequality together with $\|e_i\|_D=\sqrt{d_i[i]}$. Introduce the shorthand  \[
     \delta_{ij}:=\frac{\abs{\ccr[i]}+\gamma\sqrt{d_i[i]}}{\abs{\ccr[j]}-\gamma\sqrt{d_j[j]}} >0,
    \]
    where positivity follows from Assumption~\ref{assump:feas}.
     Set $t_j=\mathrm{sign}(\ccr[j])\abs{t_i-\hat{x}[i]}\delta_{ij}+\hat{x}[j]$. Substituting this choice into the lower bound above yields
    \begin{align*}
        h(t_i,t_j)&\geq -\gamma\abs{t_i-\hat{x}[i]}\sqrt{d_i[i]}+(t_i-\hat{x}[i])\ccr[i]+\delta_{ij}\abs{t_i-\hat{x}[i]}(\abs{\ccr[j]}-\gamma\sqrt{d_j[j]})\\
        &\geq -(\abs{\ccr[i]}+\gamma\sqrt{d_i[i]})\abs{t_i-\hat{x}[i]}+\delta_{ij}\abs{t_i-\hat{x}[i]}(\abs{\ccr[j]}-\gamma\sqrt{d_j[j]})=0.
    \end{align*}
    Thus this choice ensures $h(t_i,t_j)\geq 0$. We may define a 1-dimensional restricted objective as
    \begin{align*}
        f(t)&=\norm{g_{ij}(\hat{x})+te_i+(\mathrm{sign}(\ccr[j])\abs{t-\hat{x}[i]}\delta_{ij}+\hat{x}[j])e_j}_D^2\\
        &\qquad +\beta \norm{g_{ij}(\hat{x})+te_i+(\mathrm{sign}(\ccr[j])\abs{t-\hat{x}[i]}\delta_{ij}+\hat{x}[j])e_j}_0\\
        &=\norm{\hat{x}+(t-\hat{x}[i])e_i+\mathrm{sign}(\ccr[j])\abs{t-\hat{x}[i]}\delta_{ij}e_j}_D^2\\
        &\qquad+\beta \norm{\hat{x}+(t-\hat{x}[i])e_i+\mathrm{sign}(\ccr[j])\abs{t-\hat{x}[i]}\delta_{ij}e_j}_0.
    \end{align*}
    By the choice of $t_j$ above, the point $\hat{x}-\hat{x}[i]e_i+\mathrm{sign}(\ccr[j])\abs{\hat{x}[i]}\delta_{ij}e_j$ corresponding to $t=0$ is feasible ($h\geq 0$), while $f(\hat{x}[i])=\cF_\beta(\hat{x})$. Global optimality of $\hat{x}$ therefore implies $f(0)\geq f(\hat{x}[i])$, hence we have
    \begin{align*}
        \norm{\hat{x}}_D^2+\beta\norm{\hat{x}}_0
        &\leq \norm{\hat{x}-\hat{x}[i]e_i+\mathrm{sign}(\ccr[j])\abs{\hat{x}[i]}\delta_{ij}e_j}_D^2\\
&\quad+\beta\norm{\hat{x}-\hat{x}[i]e_i+ \mathrm{sign}(\ccr[j])\abs{\hat{x}[i]}\delta_{ij}e_j}_0.
    \end{align*}
    The vector $\hat{x}-\hat{x}[i]e_i+\mathrm{sign}(\ccr[j])\abs{\hat{x}[i]}\delta_{ij}e_j$ has at most $\abs{\hat{\sigma}}-1$ nonzero elements since the perturbation removes index $i$ without adding a new nonzero index.  Hence its $\ell_0$-term is at most $\|\hat{x}\|_0-1$, yielding
\begin{align}\label{eq:betabound}
   \notag \beta &\leq \norm{\hat{x}-\hat{x}[i]e_i+\mathrm{sign}(\ccr[j])\abs{\hat{x}[i]}\delta_{ij}e_j}_D^2-\norm{\hat{x}}_D^2\\&\notag=\tpose{(2\hat{x}-\hat{x}[i]e_i+\mathrm{sign}(\ccr[j])\abs{\hat{x}[i]}\delta_{ij}e_j)}D(-\hat{x}[i]e_i+\mathrm{sign}(\ccr[j])\abs{\hat{x}[i]}\delta_{ij}e_j)\\&
    \leq \norm{2\hat{x}-\hat{x}[i]e_i+\mathrm{sign}(\ccr[j])\abs{\hat{x}[i]}\delta_{ij}e_j}_D\norm{-\hat{x}[i]e_i+\mathrm{sign}(\ccr[j])\abs{\hat{x}[i]}\delta_{ij}e_j}_D\\
    &\notag \leq\left(2\norm{\hat{x}}_D+\abs{\hat{x}[i]}\norm{e_i-\mathrm{sign}(\hat{x}[i]\ccr[j])\delta_{ij}e_j}_D\right)\abs{\hat{x}[i]}\norm{e_i-\mathrm{sign}(\hat{x}[i]\ccr[j])\delta_{ij}e_j}_D.
\end{align}
Here, the last two inequalities follow from the Cauchy-Schwarz and the triangle inequalities. By the definition of $\rho_i$, we have $ \|e_i-\mathrm{sign}(\hat{x}[i]\ccr[j])\delta_{ij}e_j\|_D\leq \rho_i$. Substituting this together with $\|\hat x\|_D\leq \sqrt{\eta}$, established in the proof of Proposition~\ref{prop:globup}, into~\eqref{eq:betabound}, we obtain
 \begin{align}
 \label{eq:betaprop}
     \beta\leq (2\sqrt{\eta}+\abs{\hat{x}[i]}\rho_i)\abs{\hat{x}[i]}\rho_i.
 \end{align}
 Consider the polynomial $P(v)= \rho_i^2v^2+2\sqrt{\eta}\rho_i v-\beta$. Since $P$ is a quadratic with a positive leading coefficient, \eqref{eq:betaprop} implies that $|\hat x[i]|$ must be larger than the positive root of $P$. Evaluating these roots gives
 \begin{align*}
     v_+,v_-=\frac{-2\sqrt{\eta}\rho_i\pm\sqrt{4\eta\rho_i^2+4\rho_i^2\beta}}{2\rho_i^2}=\frac{-\sqrt{\eta}\pm\sqrt{\eta+\beta}}{\rho_i}.
 \end{align*}
 Hence, we obtain the lower bound $\abs{\hat{x}[i]}\geq \frac{\sqrt{\eta+\beta}-\sqrt{\eta}}{\rho_i}$.

 \textbf{Case 2: } Assume $\abs{\hat{\sigma}}=1$, then we have a closed form solution for the specific entry $i\in\hat{\sigma}$, which takes the form
\begin{align*}
    \hat{x}[i]=\frac{\bar{\ccr}\;\mathrm{sign}(\ccr[i])}{\abs{\ccr[i]}-\gamma\sqrt{d_i[i]}}\quad\Rightarrow\quad \abs{\hat{x}[i]}=\frac{\bar{\ccr}}{\abs{\ccr[i]}-\gamma\sqrt{d_i[i]}}.
\end{align*}
In both cases the lower bound takes the form
\[
    \min\left\{ \frac{\sqrt{\eta+\beta}-\sqrt{\eta}}{\rho_i},\frac{\bar{\ccr}}{\abs{\ccr[i]}-\gamma\sqrt{d_i[i]}}\right\},
\] where the first term is active when $\abs{\hat{\sigma}}\geq2$ and the second when $\abs{\hat{\sigma}}=1$. Since a global minimizer must fall into one of these two cases, the bound holds unconditionally for every $i\in\hat{\sigma}$.
 \end{proof}

\end{document}
